\section{Evidence base and closest predecessors}\label{sec:evidence}

\subsection{Scope and verification}

\noindent The literature was searched on November~9--10, 2026, extending earlier searches with targeted queries on contest
competition, coalitional conflict, leveling, the validity of physical attractiveness cues, the antiquity of games, late
hominins, and selection in contemporary populations. This is an integrative review, not a systematic one; no pooled effect is
estimated. Bibliographic records for all \RefResolved{} references with a DOI were retrieved from Crossref and the reference
list is generated from them. Specific claims were checked against the abstracts of the cited sources where an abstract could
be retrieved (Section~\ref{app:repro}).

\subsection{Preferences, perception, and validity}

Intelligence ranks among the most desired partner traits for both sexes across cultures, alongside average sex differences in
other preferences~\cite{buss1989,walter2020}. Evidence linking measured intelligence to immediate appeal is mixed. A small study
found intelligent men consistently appealing~\cite{prokosch2009}, but larger video and live-interaction studies found that
measured intelligence was detectable yet did not predict appeal, whereas perceived intelligence and funniness
did~\cite{driebe2021}. Overconfidence can raise desirability while risking perceived arrogance~\cite{murphy2015}. Humor
production is predicted by intelligence and predicts lifetime sexual partners~\cite{greengross2011}; musical aptitude showed no
positive association with partner number in more than 10,000 Swedish twins~\cite{mosing2015}. Stated ideals predict
in-person desire poorly~\cite{eastwick2008}. Two results bear on the central premise. Judgments of intelligence from short
behavioral episodes predict test scores, and prediction improves as more episodes are pooled, with small gains beyond about
six~\cite{borkenau2004}: the concave return to repetition assumed in the model. And ratings by others predict a person's
outcomes better than self-reports, with accuracy that rises with acquaintance~\cite{connelly2010}: the return to long
observation. The lesson is
that validity is a parameter, not an assumption. Displays can be persuasive without being informative.

\subsection{Performance, standing, and reproduction in small-scale societies}

Table~\ref{tab:field} collects the field evidence most relevant to the mechanism. In a meta-analysis of 288 results from 33
nonindustrial populations, men's status predicted reproductive success with a pooled correlation of about 0.19; effects were
similar whether status rested on hunting ability, political influence, wealth, or formidability, and similar across
subsistence types~\cite{vonrueden2016}. Better hunters have higher age-corrected reproductive success in all five societies
reviewed by Smith~\cite{smith2004}; whether men hunt mainly to provision their families or to display is
debated~\cite{hawkes2002,gurvenhill2009}, and the mechanism here needs only that hunting be observed and valued. Among Tsimane forager-horticulturalists both dominant and prestigious men have more
surviving offspring, and the strongest path from prestige to fitness is marriage to women who begin reproducing
earlier~\cite{vonrueden2011}, which is the Darwin--Fisher route of the model below. Among Agta foragers, skilled storytellers
were nearly twice as likely to be named as preferred campmates and had 0.53 more living offspring~\cite{smith2017}. Sexual
selection operates even under imposed monogamy, as a historical Finnish population shows~\cite{courtiol2012}, while a
meta-analysis finds that masculine traits predict mating more consistently than reproduction~\cite{lidborg2022}. On the
women's side, a Tsimane' mother's plant knowledge predicts her children's health independently of schooling and
wealth~\cite{mcdade2007}, a direct return to a woman's own capacity. These are correlational results. None measures cognitive
capacity independently of standing, which is the key missing parameter (Section~\ref{sec:strength}).

\begin{table*}[t]
\caption{Field and population evidence linking observable performance or standing to reproduction.}\label{tab:field}
\scriptsize
\begin{tabularx}{\textwidth}{@{}P{0.155\textwidth}P{0.15\textwidth}P{0.175\textwidth}LP{0.20\textwidth}@{}}
\toprule
Study & Population & Performance or standing & Reproductive association & Principal limitation\\
\midrule
von Rueden \& Jaeggi~\cite{vonrueden2016} & 33 nonindustrial populations & Status: hunting, influence, wealth, formidability &
Pooled $r\approx0.19$ (288 results); similar across status types & Men only; cognition not measured\\
Smith~\cite{smith2004} & Ache, Hadza, !Kung, Lamalera, Meriam & Hunting success & Better hunters have higher age-corrected
reproductive success & Provisioning, signaling and phenotypic correlation not separated\\
Smith, Bliege Bird \& Bird~\cite{smith2003} & Meriam, Torres Strait & Public, costly turtle hunting & Earlier first
reproduction, more and higher-quality mates & No manipulation\\
von Rueden et al.~\cite{vonrueden2011} & Tsimane, Bolivia & Dominance and prestige (peer rated) & Both raise in-pair surviving
offspring; prestige via wives who reproduce earlier & One society; status causes and consequences entangled\\
Apicella~\cite{apicella2014} & Hadza, Tanzania & Upper-body strength, hunting reputation & Stronger men have more children &
Strength, not cognition\\
Smith et al.~\cite{smith2017} & Agta, Philippines & Storytelling skill (peer nominations) & 0.53 more living offspring;
preferred campmates & Correlational\\
Llaurens et al.~\cite{llaurens2009} & Sereer, Senegal & Traditional wrestling & More offspring among participants after
covariates & Prestige and family decisions; injury costs\\
Glowacki \& Wrangham~\cite{glowacki2015} & Nyangatom, Ethiopia & Raiding record & Elders who raided more have more wives and
children & Works through bridewealth; pastoralists\\
Beckerman et al.~\cite{beckerman2009} & Waorani, Ecuador & Raiding record & More aggressive warriors have \emph{lower}
reproductive success & Extreme homicide rates\\
Gurven et al.~\cite{gurven2014} & Tsimane, Bolivia & Personality factors & Several factors raise men's fertility;
industriousness has an intermediate optimum & Self-report; varies by region in women\\
McDade et al.~\cite{mcdade2007} & Tsimane', Bolivia & Mother's plant knowledge & Better child growth and immune markers &
Offspring quality, not number\\
Greengross \& Miller~\cite{greengross2011} & U.S.\ students & Humor production & More lifetime sexual partners & Mating, not
fertility\\
Mosing et al.~\cite{mosing2015} & Swedish twins & Musical aptitude & No positive association with partner number &
Low-fertility modern population\\
\bottomrule
\end{tabularx}
\end{table*}

\subsection{Assessment is multi-party, long-windowed, and pooled by talk}

Across 190 hunter-gatherer societies, parents or close kin arrange marriages in about 85\%, and brideservice, brideprice, or
other exchange occurs in about 80\%; phylogenetic reconstruction infers arranged marriage at least since the out-of-Africa
dispersal~\cite{walker2011}. Parental choice is itself a selective force on mating~\cite{apostolou2007}. This weakens any model
that assumes autonomous mate choice, and it strengthens the information argument: kin evaluate a prospective son- or
daughter-in-law over years, and brideservice places a young man's provisioning under the direct observation of his in-laws
before the marriage is secure. Prestige, freely conferred on those whose skill others wish to learn from~\cite{henrich2001},
is a parallel assessment market.

The audience is large. Forager bands are small, but they are nodes in networks of hundreds: Hadza and Ache men are
estimated to watch over 300 other men make tools in a lifetime, against about 20 male partners for a male
chimpanzee~\cite{hill2014}, and co-residents are mostly unrelated~\cite{hill2011}. Information also travels secondhand. About
two-thirds of conversation time is given to social topics~\cite{dunbar2004}, and gossip changes how people treat a third
party even when they have observed that party directly~\cite{sommerfeld2007}. Language multiplies the observers of every
performance at almost no cost, and it also multiplies shared error, a trade-off formalized in Section~\ref{sec:portfolio}.

\subsection{Venues for performance}

The ancestral analogue of games is a performance ecology: foraging returns that are publicly shared and discussed,
tool-making, tracking, play, contests, song, dance, and talk. Children in foraging societies learn subsistence skills in
mixed-age play groups, and adults begin to teach hunting only in adolescence~\cite{lewlevy2017}, so a decade of public
practice precedes marriage. Fire extended the social day. Among Ju/'hoansi foragers daytime talk centered on work and gossip,
whereas over 80\% of firelit talk concerned stories, song, dance, and spiritual matters~\cite{wiessner2014}. Fire use is
documented at Wonderwerk Cave about one million years ago~\cite{berna2012}, and habitual use in Europe from roughly
400,000--300,000 years ago~\cite{roebroeks2011}. Adult playfulness has been studied as a mate
signal~\cite{chick2012,moraes2022}, sport participation shifts perceived mate characteristics~\cite{schultehostedde2012}, mating
games and sport have been analyzed as culturally evolving signals~\cite{deblock2007,deblock2009}, and quick verbal wit has
been proposed as a target of selection since early grammar~\cite{progovac2026}. Games of strategy track political complexity
across cultures, while games of physical skill are widespread~\cite{roberts1959}.

\subsection{Contest, coalition, and leveling}

Three literatures bear on the narrative of Section~\ref{sec:audit}. First, animal contest theory explains why fights are
usually restrained: contestants assess one another through graded, repeated displays and escalate only when closely matched or
when the future is worth little~\cite{maynardsmith1973,parker1974,enquist1983,enquist1990,arnott2009}. Second, human men bear
the marks of contest competition, in muscle mass, in the accuracy with which strength is read from bodies and faces, and in
the link between dominance and mating success~\cite{puts2010,sell2009,lassek2009,hill2013}; yet human sexual dimorphism in body
and canine size is modest among primates~\cite{plavcan2012}, and the species is better described as one of mutual mate choice
with moderate sex differences than as a contest specialist~\cite{stewartwilliams2013}. Third, human groups level their
dominants. Foragers ridicule, desert and occasionally kill would-be despots~\cite{boehm1993,boehm1999}; models show that
coalitions of the weak can evolve to curb the strong~\cite{gavrilets2012,pandit2003}, that reproductive leveling favors
group-beneficial traits under intergroup competition~\cite{bowles2006}, and that within-group conflict falls as between-group
conflict rises~\cite{reeve2007}. Across primates male mating skew falls as the number of males in a group
rises~\cite{kutsukake2006}, and human reproductive skew is low for a mammal~\cite{ross2023}. How often Pleistocene groups
fought is disputed: estimates of mortality from intergroup violence that are high enough to drive the evolution of
altruism~\cite{bowles2009} sit beside evidence that most lethal events among mobile foragers are personal
disputes~\cite{fry2013}, a phylogenetic baseline of about 2\% of deaths from conspecific violence~\cite{gomez2016}, a massacre
of foragers about 10,000 years ago~\cite{lahr2016}, and lethal interpersonal violence 430,000 years ago~\cite{sala2015}.

\subsection{Comparative evidence}

Mate choice for cognitive traits has been reviewed across nonhuman vertebrates~\cite{boogert2011}. Female budgerigars shifted
preference toward males seen solving problems~\cite{chen2019}. In budgerigars a composite score across four cognitive tasks
predicted male reproductive output, whereas individual tasks mostly did not~\cite{medinagarcia2021}, the pattern the portfolio
theorem predicts. In mosquitofish, paternity analysis of 2,430 offspring showed that inhibitory control and spatial learning
predicted siring success while other cognitive traits did not~\cite{vinogradov2025}. A zebra-finch study found that dance
complexity tracked body condition and attractiveness, not cognitive performance~\cite{baroudagues2025}: displays can carry
condition rather than capacity.

Phylogenetic comparisons cut against any simple claim that sexual selection enlarges brains. Across primates, relative
brain size is unrelated to sperm competition and \emph{negatively} related to male competition for mates, and the largest
relative brains occur in monogamous species~\cite{schillaci2006}. Across mammals, brain mass does not vary with the level of
sperm competition~\cite{lemaitre2009}; in pinnipeds, stronger sexual selection goes with smaller relative brains in males
because bodies enlarge, not because brains shrink~\cite{fitzpatrick2012}. Outside primates, larger relative brains go with
pair-bonded mating systems~\cite{shultz2007}. Experiments supply two further facts. Artificial selection on relative brain
size in guppies produced a rapid response, better learning in large-brained females, smaller guts, and fewer
offspring~\cite{kotrschal2013}: brain tissue is heritable, functional, and costly in a vertebrate. And large-brained female
guppies discriminate male quality better~\cite{corrallopez2017}: assessment is itself cognitively demanding.

\subsection{Rival mechanisms that produce the same appeal}

Three mechanisms can generate the attraction patterns the hypothesis predicts without valid information about capacity.
Display and preference can coevolve by runaway, so a chosen display need not reveal anything~\cite{lande1981}. Observers can
copy others' choices, inferring attractiveness from popularity; models show this alone can maintain variation in traits and
preferences~\cite{duval2023}, and a meta-analysis confirms that humans copy mate choices~\cite{gouda2018}. And status can
raise both training access and mating success. These rivals predict gains in appeal without gains in accuracy, which is why
the decisive experiment crosses information with endorsement (Section~\ref{sec:predictions}).

\subsection{Energetics, development, and genetics}

Humans have elevated total energy expenditure relative to other apes~\cite{pontzer2016}; cooking and processing raise energy
returns and reduce chewing~\cite{carmody2011,zink2016}; comparative data do not support a simple brain--gut
exchange~\cite{navarrete2011}; and cooperative care and grandmothering spread the cost of slow
development~\cite{hawkes1998,burkart2014}. The mother's body is part of the budget. Human infants are unusually fat, which
buffers the brain's energy supply~\cite{kuzawa1998,cunnane2003}; gestation length and neonatal size appear limited by maternal
metabolism as much as by the pelvis~\cite{dunsworth2012}; and pelvis shape, stature and head size covary so that women with
large heads, who tend to bear large-headed infants, have birth canals that accommodate them~\cite{fischer2015}, under selection
that is weak on average but severe at the margin~\cite{mitteroecker2016}. Prolonged brain growth may already have been present
in \textit{Australopithecus afarensis}~\cite{gunz2020} and extended growth in early \Homo~\cite{zollikofer2024}. Brain shape and
size have distinct histories~\cite{neubauer2018}; a human-like organization of the frontal lobes appears only 1.7--1.5~million
years ago, after the first dispersal from Africa~\cite{poncedeleon2021}; and developmental dynamics can generate enlargement
without direct selection on volume~\cite{gonzalezforero2024}.

Within contemporary humans, brain volume correlates with intelligence at about $r=0.24$ in meta-analysis~\cite{pietschnig2015}
and $r=0.19$ in a preregistered sample of 13,608~\cite{nave2019}; the genetic correlation is about
0.24~\cite{jansen2020}. Such genetic correlations can be inflated by cross-trait assortative mating~\cite{border2022}, but the
association persists within families at about the same size, the pattern expected if larger brains contribute causally to
ability~\cite{lee2019}. Partner correlations are about 0.44 for IQ score and 0.55 for educational
attainment~\cite{horwitz2023}, and assortment leaves genomic signatures~\cite{robinson2017}.

\subsection{Closest theoretical predecessors}

Mutual choice for cognitive indicators~\cite{miller2000ch,miller2000book}, costly signaling under
monogamy~\cite{hooper2008,grafen1990}, mutual choice under direct benefits~\cite{courtiol2016}, sexual selection through
breeding date in monogamous species~\cite{kirkpatrick1990}, social selection by partner
choice~\cite{westeberhard1983,nesse2007}, and runaway social competition~\cite{alexander1989,flinn2005,crespi2022} are all
established. Contest as the main mode of sexual selection in men~\cite{puts2010}, leveling~\cite{boehm1999,gavrilets2012},
and Lanchester's laws in human combat~\cite{johnson2015} are each developed elsewhere. Mating games and sport as culturally
evolving signals have been analyzed directly~\cite{deblock2007,deblock2009}, strategic and assessment-driven models of
intelligence exist~\cite{gavrilets2006,baroudagues2022}, and the mathematics of response to selection on an imperfect index is
eighty years old~\cite{lush1937,hazel1943,crow1958}. What is added here is a quantitative account of how information limits
and aims social selection, how it accumulates across generations, how contest and coalition change its intensity and its
target, and how it interacts with the energetic budget to produce persistent, stepwise investment in tissue.
